\documentclass[a4wide,11pt]{report}
\usepackage[utf8]{inputenc}
\usepackage{graphicx}
\usepackage{url}
\usepackage{float}
\usepackage{csvsimple}
\usepackage[pdfborder=0 0 0]{hyperref}
\usepackage{longtable}
\usepackage{fancyhdr}
\pagestyle{fancy}
\usepackage{graphicx}

\begin{document}

\title{Standard Grimoire Report \\
  TemplateProject Project \\
      2016-Q3 \\
        ~~\\~~\\  \includegraphics[scale=.35]{logo}       \\}




\maketitle


\newpage
~~\\~\\~\\~\\
%\begin{tabular}{p{5cm}p{10cm}}
%~  & Daniel Izquierdo Cort\'azar \\
%       & Chief Data Officer  \\
%       & dizquierdo@bitergia.com \\

%\end{tabular}

~~\\~~\\

~~\\
~~\\

\begin{tabular}{p{5cm}p{10cm}}

\end{tabular}

~~\\
~~\\

This report would not exist without the effort of the people involved in the development of the \href{http://grimoirelab.github.io}{Grimoire toolset}.

~~\\
~~\\
~~\\
~~\\
~~\\
~~\\

\begin{flushright}
  (cc) 2016 Bitergia. Some rights reserved. \\
  This work licensed under Creative Commons
  Attribution-ShareAlike 4.0 Unported License. \\
  To view a copy of full license, see
  \url{http://creativecommons.org/licenses/by-sa/4.0}, \\
  or write to
  Creative Commons, 559 Nathan Abbott Way, Stanford,
  California 94305, USA. \\ \ 
\end{flushright}


%\rhead{Bitergia Report}
%\rhead{\includegraphics[scale=.15]{openstack.eps}\includegraphics[scale=.15]{logo.eps}}
\rhead{\includegraphics[scale=.15]{logo}}
%\rhead{\includegraphics[scale=.15]{logo.eps}}


\newpage

\begin{center}
\textbf{Executive Summary}
\end{center}

This report provides a quantitative analysis of the current and past situation of the TemplateProject project. All the data presented in it is based on information retrieved from the software development repositories of the project. The analysis includes a summary of the general situation of the project, and specific analysis of some of its development processes (Gerrit) and communication channels (mailing lists). For comparison with the past, most of the data is shown on a comparison with previous periods.

\setcounter{tocdepth}{4}
\setcounter{secnumdepth}{4}
\tableofcontents


\chapter{Project overview}

The report looks at activities across the TemplateProject community during
2016-Q3, comparing it to previous
period of analysis.

\begin{table}[H]
    \centering
    \begin{tabular}{c|c|c|}%
    \bfseries Data source & \bfseries Activity last quarter & \bfseries Change (wrt to prev. quarter) % specify table head
    \csvreader[head to column names]{data/data_source_evolution.csv}{}% use head of csv as column names
    {\\\datasource & \netvalues ~ \metricsnames & \relativevalues\% }
    \end{tabular}
    \caption{Activity during the last period of analysis and its evolution}
    \label{tab:_summary}
\end{table}

Table~\ref{tab:_summary} shows development activity for each of the analyzed
data sources. The activity column displays information about the net
activity numbers, while the Change column displays information
about the relative difference with respect to the previous period of
analysis.

%TODO: add list of repositories analyzed per quarter

%\section{Development and Conversations}

%The mean number of developers per month

%\begin{figure}[H]
%    \centering
%    \includegraphics[scale=.35]{figs/authors_month.eps}
%\caption{Number of developers per month in mean for each quarter}
%\end{figure}

\begin{tabular}{p{7cm} p{5cm}}
    \vspace{0pt} 
    \includegraphics[scale=.35]{figs/git_authors.eps}
    & 
    \vspace{0pt}
    \begin{tabular}{l|r|}%
    \bfseries Period & \bfseries Authors per quarter% specify table head
   \csvreader[head to column names]{data/git_authors.csv}{}% use head of csv as column names
   {\\\labels & \authors}
   \end{tabular}
  
\end{tabular}




%The total number of contributors divided into three sets (core,
%regular and casual\footnote{Contributing developers are characterized
%as core, regular and casual depending on their activity in the git
%repositories. The classification is built by sorting contributors by
%their total number of commits; we sum the total commits per each
%individual contributors: the individuals whose commits sum up to 80\%
%of the total number of commits in the quarter are the core
%contributors in that quarter. The regular contributors are those whose
%commits sum up to 95\% of the total. The others are the casual
%contributors.}) follow a similar pattern.

%\begin{figure}[H]
%    \centering
%    \includegraphics[scale=.35]{figs/onion.eps}
%    \caption{Evolution during the last quarters of core, regular and casual developers (based on git activity)}
%\end{figure}

%\begin{table}[H]
%    \centering
%    \begin{tabular}{l|r|r|r|}%
%    \bfseries Period & \bfseries Core & \bfseries Regular & \bfseries Occasional% specify table head
%    \csvreader[head to column names]{data/onion_model.csv}{}% use head of csv as column names
%    {\\\labels & \core & \regular & \occasional}
%    \end{tabular}
%    \caption{Characterization of developers by their total contribution to the project}
%\end{table}



This report aims to provide some insight into the software
development process of the TemplateProject community measuring efficiency
and process of the community based on three metrics: the Review
Efficiency Index (REI), the Time to Merge (TTM). 
REI is measured as the number of closed
changesets out of the new ones in a
given period. TTM is measured as the time since a review is submitted
until this is closed. 


\begin{table}[H]
    \centering
    \begin{tabular}{c|c|}%
    \bfseries  REI  & \bfseries TTM % specify table head
    \csvreader[head to column names]{data/efficiency.csv}{}% use head of csv as column names
    {\\\bmireviews  & \daystomergereviewmedian ~ days }
    \end{tabular}
    \caption{Closed changesets out of opened changesets (REI) and median time to merge (TTM)}
\end{table}

%\chapter{Communication and support-related activities}

%Analysis of the communication channels used for communication and support-related activities.


The following charts show activity in terms of emails sent and people sending emails. In addition, a table is presented with the hot topics in the several analyzed mailing lists. This shows hot topics ordered by number of total posts in such thread.

\begin{tabular}{p{7cm} p{5cm}}
    \vspace{0pt} 
    \includegraphics[scale=.35]{figs/mls_emails_sent.eps}
    & 
    \vspace{0pt}
    \begin{tabular}{l|l}%
    \bfseries Period & \bfseries Emails % specify table head
    \csvreader[head to column names]{data/mls_emails_sent.csv}{}% use head of csv as column names
    {\\\labels & \emailssent}
    \end{tabular}
\end{tabular}

\begin{tabular}{p{7cm} p{5cm}}
    \vspace{0pt} 
    \includegraphics[scale=.35]{figs/mls_emails_senders.eps}
    & 
    \vspace{0pt}
    \begin{tabular}{l|l}%
    \bfseries Period & \bfseries People % specify table head
    \csvreader[head to column names]{data/mls_emails_senders.csv}{}% use head of csv as column names
    {\\\labels & \emailssenders}
    \end{tabular}
\end{tabular}

%\begin{tabular}{p{7cm} p{5cm}}
%    \vspace{0pt} 
%    \includegraphics[scale=.35]{figs/emails_senders_init.eps}
%    & 
%    \vspace{0pt}
%    \begin{tabular}{l|l}%
%    \bfseries Period & \bfseries People % specify table head
%    \csvreader[head to column names]{data/emails_senders_init.csv}{}% use head of csv as column names
%    {\\\labels & \senders}
%    \end{tabular}
%\end{tabular}

%\noindent{\footnotesize
%\begin{tabular}{p{2.5cm}p{7.5cm}p{1.5cm}}
%    \bfseries Initial Author &  & \bfseries Number \\ 
%    \bfseries and Date       & \bfseries Subject  & \bfseries Messages% specify table head
%    \csvreader[head to column names]{data/mls_top_longest_threads.csv}{}% use head of csv as column names
%    {\\\initiator \\\date & \subject & \len}
%   % {\\\hline\csvcoli & \csvcolvii & \csvcolv & \csvcolii}% specify your coloumns here
%\end{tabular}
%}

%% \noindent{\footnotesize
%% \begin{tabular}{p{3cm}p{6cm}p{2cm}}
%%     \bfseries Initial Author & & \bfseries Diff. \\
%%     \bfseries and Date & \bfseries Subject  &  \bfseries People% specify table head
%%     \csvreader[head to column names]{data/mls_top_crowded_threads.csv}{}% use head of csv as column names
%%     {\\\initiator \\\date & \subject  & \people}
%%    % {\\\hline\csvcoli & \csvcolvii & \csvcolv & \csvcolii}% specify your coloumns here
%% \end{tabular}
%% }

%\begin{tabular}{p{5cm}p{3cm}p{2cm}}
%    \bfseries Initial Author & \bfseries Subject  & \bfseries Date % specify table head
%    \csvreader[head to column names]{data/mls_top_longest_threads.csv}{}% use head of csv as column names
%    {\\\hline\csvcoli & \csvcolvi & \csvcoliv }% specify your coloumns here
%\end{tabular}


The rest of the document is divided into three chapters with information from the last periods: 

\begin{itemize}
\item activity: centered on the following metrics: commits from git activity, submitted and closed changests from the review system. 
\item community: active authors in git and top ten developers and top ten organizations contributing to the development in the last period of analysis.
\item process: efficiency closing changesets and Time to Merge (mean and median).
\end{itemize}

\chapter{Activity}
Commits in Git, submitted and closed reviews.


\begin{tabular}{p{7cm} p{5cm}}
    \vspace{0pt} 
    \includegraphics[scale=.35]{figs/git_commits_gerrit_submitted_general.eps}
    & 
    \vspace{0pt}
    \begin{tabular}{l|r|r|}%
    \bfseries Period & \bfseries Commits & \bfseries Reviews% specify table head
    \csvreader[head to column names]{data/git_commits_gerrit_submitted_general.csv}{}% use head of csv as column names
    {\\\labels & \commits & \submitted}
    \end{tabular}
\end{tabular}


%\begin{tabular}{p{7cm} p{5cm}}
%    \vspace{0pt} 
%    \includegraphics[scale=.35]{figs/its_opened_its_closed_general.eps}
%    & 
%    \vspace{0pt}
%    \begin{tabular}{l|r|r|}%
%    \bfseries Period & \bfseries Closed & \bfseries Opened% specify table head
%    \csvreader[head to column names]{data/its_opened_its_closed_general.csv}{}% use head of csv as column names
%    {\\\labels & \closed & \opened}
%    \end{tabular}
%\end{tabular}

\begin{tabular}{p{7cm} p{5cm}}
    \vspace{0pt} 
    \includegraphics[scale=.35]{figs/gerrit_submitted_gerrit_closed_general.eps}
    & 
    \vspace{0pt}
    \begin{tabular}{l|r|r|}%
    \bfseries Period & \bfseries Opened & \bfseries Closed % specify table head
    \csvreader[head to column names]{data/gerrit_submitted_gerrit_closed_general.csv}{}% use head of csv as column names
    {\\\labels & \submitted & \closed}
    \end{tabular}
\end{tabular}

\chapter{Community}
Active authors in Git, top authors and organizations in the last period of analysis.


\begin{tabular}{p{7cm} p{5cm}}
    \vspace{0pt} 
    \includegraphics[scale=.35]{figs/git_authors_general.eps}
    & 
    \vspace{0pt}
    \begin{tabular}{l|l}%
    \bfseries Period & \bfseries Active Authors % specify table head
    \csvreader[head to column names]{data/git_authors_general.csv}{}% use head of csv as column names
    {\\\labels & \authors}
    \end{tabular}
\end{tabular}


\begin{tabular}{p{7cm} p{5cm}}
    \vspace{0pt}
\begin{tabular}{l|l}%
    \bfseries Author & \bfseries Commit (s) % specify table head
    \csvreader[head to column names]{data/git_top_authors_general.csv}{}% use head of csv as column names
    {\\\hline\csvcoli&\csvcolii}% specify your coloumns here
\end{tabular}
\end{tabular}

\begin{tabular}{p{7cm} p{5cm}}
\vspace{0pt}
\begin{tabular}{l|l}%
    \bfseries Organization & \bfseries Commit (s) % specify table head
    \csvreader[head to column names]{data/git_top_organizations_general.csv}{}% use head of csv as column names
    {\\\hline\csvcoli&\csvcolii}% specify your coloumns here
\end{tabular}
\end {tabular}


\chapter{Process}
Efficiency closing changesets and time to review (mean and median).


%\begin{tabular}{p{7cm} p{5cm}}
%    \vspace{0pt} 
%    \includegraphics[scale=.35]{figs/its_bmi_tickets_general.eps}
%    & 
%    \vspace{0pt}
%    \begin{tabular}{l|l}%
%    \bfseries Period & \bfseries Closed/Opened % specify table head
%    \csvreader[head to column names]{data/its_bmi_tickets_general.csv}{}% use head of csv as column names
%    {\\\labels & \bmitickets}
%    \end{tabular}
%\end{tabular}

\begin{tabular}{p{7cm} p{5cm}}
    \vspace{0pt} 
    \includegraphics[scale=.35]{figs/gerrit_bmi_reviews_general.eps}
    & 
    \vspace{0pt}
    \begin{tabular}{l|l}%
    \bfseries Period & \bfseries Closed/Subm. % specify table head
    \csvreader[head to column names]{data/gerrit_bmi_reviews_general.csv}{}% use head of csv as column names
    {\\\labels & \bmireviews}
    \end{tabular}
\end{tabular}


%\begin{tabular}{p{7cm} p{5cm}}
%    \vspace{0pt} 
%    \includegraphics[scale=.35]{figs/its_days_to_close_ticket_avg_its_days_to_close_ticket_median_general.eps}
%    & 
%    \vspace{0pt}
%    \begin{tabular}{l|r|r|}%
%    \bfseries Period & \bfseries Median & \bfseries Mean % specify table head
%    \csvreader[head to column names]{data/its_days_to_close_ticket_avg_its_days_to_close_ticket_median_general.csv}{}% use head of csv as column names
%    {\\\labels & \daystocloseticketmedian & \daystocloseticketavg}
%    \end{tabular}
%\end{tabular}



\begin{tabular}{p{7cm} p{5cm}}
    \vspace{0pt} 
    \includegraphics[scale=.35]{figs/gerrit_days_to_merge_review_avg_gerrit_days_to_merge_review_median_general.eps}
    & 
    \vspace{0pt}
    \begin{tabular}{l|r|r|}%
    \bfseries Period & \bfseries Median & \bfseries Mean % specify table head
    \csvreader[head to column names]{data/gerrit_days_to_merge_review_avg_gerrit_days_to_merge_review_median_general.csv}{}% use head of csv as column names
    {\\\labels & \daystomergereviewmedian & \daystomergereviewavg}
    \end{tabular}
\end{tabular}





\appendix{}

\chapter{Metrics Definitions}
\label{chap:metrics_definitions}

\begin{itemize}
\item Commit: this is defined as the action(s) that performes a change in the source code. 
Bots, merges and other type of automatic activity is removed from the records. In addition,
when aggregating several git repositories, this metric only counts unique revisions (unique hashes found in the git repositories).
In addition, all branches are aggregated to the analysis.

\item Submitted changesets: a code review is the process of peer reviewing source code changes. A submitted code
is not merged to the master code of a given project till this is approved.
A submitted code review is defined as any changeset submitted to the Gerrit
system. 
 
\item Authors: a developer is defined as author if she is the owner of the patchset sent for reviewing and this is 
merged into the source code. As previously indicated, automatic commits such bot's are removed from this analysis.

\item Efficiency closing changesets: this metric is a derivation of the Backlog Management Index as it is named as
Review efficiency index (REI). As similarly used in the BMI index, this metrics measures the number of closed changesets
 out of the total number of new changesets in a given period.

\item Time to Merge: this time consists of the time between the first upload of the first changeset
till the last iteration of the code review process is merged into the code.
This metric is provided in number of days.


\end{itemize}


Metrics measured in the general overview:

\begin{itemize}
\item Developer per period: average of developers per period ignoring bots and merges.

\item Emails sent: number of emails sent by people to the several mailing lists. Bots are not registered.

\item People sending emails: number of people sending those emails ignoring bots.

\end{itemize}


%\chapter{Source code and data sources}
%\label{chap:data_sources}

%The source code of the scripts and templates used to produce this
%report are available from the GrimoireReports repository\footnote{\url{
%https://github.com/grimoirelab}}.
%Commit identified with the hash \emph{d33b6f5951662fb5bd895eee89c0b0344b09c956}. 

%select p.id, pr.data_source, pr.repository_name 
%from projects p, project_repositories pr 
%where p.project_id=pr.project_id order by p.id, pr.data_source;

%\input{repositories.tex}

\end{document}
